\documentclass[11pt,a4paper]{article}
\usepackage[T1]{fontenc}
\usepackage[utf8]{inputenc}
\usepackage[margin=0.9in]{geometry}
\usepackage{enumitem}
\usepackage{titlesec}
\usepackage{booktabs}
\usepackage{tabularx}
\usepackage{array}
\usepackage{amssymb}
\usepackage[hidelinks]{hyperref}

\newcommand{\studentname}{Shawal Kabir Chy.}

\linespread{1.05}
\titlespacing*{\section}{0pt}{14pt}{6pt}
\titleformat{\section}{\normalfont\Large\bfseries}{\thesection.}{0.5em}{}
\setlist[itemize]{topsep=4pt,itemsep=3pt,parsep=0pt,leftmargin=*}
\setlist[enumerate]{topsep=4pt,itemsep=3pt,parsep=0pt,leftmargin=*}
\setlength{\parindent}{0pt}
\setlength{\parskip}{6pt}
\newcolumntype{L}{>{\raggedright\arraybackslash}X}
\pagestyle{plain}

\begin{document}

{\centering\LARGE\bfseries Project Roadmap\par}
\vspace{6pt}
{\centering\large Silent Tool-Call Failures: How Quantization Degrades Function Calling\\
and Structured Output in Small LLMs\par}
\vspace{12pt}
{\centering
\begin{tabular}{rl}
\textbf{Course:}     & CSE498R Directed Research \\
\textbf{Supervisor:} & Dr.\ Mohammad Shifat-E-Rabbi \\
\textbf{Student:}    & \studentname \\
\textbf{Repository:} & \href{https://github.com/shawalkabirchy/quantization-functional-call}{\texttt{github.com/shawalkabirchy/quantization-functional-call}} \\
\textbf{Timeline:}   & 4 October -- 2 December 2026 (9 weeks) \\
\end{tabular}\par}
\vspace{6pt}

% -------------------------------------------------------------------
\section{Overview}

The project runs for nine weeks, from 4 October to 2 December 2026, and ends with a complete paper
and a public code release. Work is pushed to the repository continuously, with two formal updates
each week, presented in person in the supervisor's office every Monday and Wednesday at 4:00 pm
(Bangladesh time), 18 in total. Each update is pushed before the meeting, adds a dated entry to
\texttt{UPDATES.md} and includes at least one concrete result: a table, a figure, or working code.
New experiments stop on 22 November, leaving the final two weeks for writing and release.

A Monday update follows about five working days, including the weekend, while a Wednesday update
follows only two. New code and GPU-heavy runs are therefore scheduled before Monday meetings, and
analysis, plotting and writing before Wednesday meetings.

% -------------------------------------------------------------------
\section{Scope: Core and Stretch}

The original proposal covers nine models, five precisions, a KV-cache quantization axis, four
benchmarks, constrained decoding and a verifier. That is more than nine weeks allow, so the work is
split into a \emph{core} that the paper requires and \emph{stretch} items attempted only if time
remains.

\begin{center}
\renewcommand{\arraystretch}{1.3}
\begin{tabularx}{\textwidth}{LL}
\toprule
\textbf{Core (required for the paper)} & \textbf{Stretch (only if time allows)} \\
\midrule
BFCL v4 AST categories, plus relevance and irrelevance detection & API-Bank, ToolACE \\
Qwen3-1.7B/4B/8B (thinking off), Llama-3.2-1B/3B-Instruct & Qwen3-0.6B, Phi-4-mini, xLAM-2-1b/8b \\
FP16, INT8, NF4, GPTQ-4bit, AWQ-4bit & INT8/INT4 KV-cache quantization \\
Syntactic vs.\ semantic failure taxonomy & Log-probability + self-check verifier \\
Constrained vs.\ unconstrained decoding (headline result) & Full 400-case stress suite \\
Schema-stress suite, about 150 cases & \\
Memory and latency recorded for every run & \\
\bottomrule
\end{tabularx}
\end{center}

The core alone is 5 models $\times$ 5 precisions $\times$ 2 decoding modes $=$ 50 evaluation runs,
which fits the timeline.

% -------------------------------------------------------------------
\section{Week-1 Prerequisites}

\begin{itemize}
  \item \textbf{GPU access.} No lab GPU is available, and vLLM does not run natively on Windows.
        All experiments therefore run on Kaggle (free, about 30 GPU-hours per week, 2$\times$ T4),
        with Google Colab as backup. The 8B model in FP16 (about 16\,GB) is split across both T4s.
        T4 GPUs do not support BF16, so FP16 serves as the full-precision baseline, and the paper
        states this.
  \item \textbf{Novelty check (done 5 Oct).} A fresh search of arXiv, Google Scholar and GitHub
        found earlier work on quantization and tool calling: ACBench (ICML 2025), Less is More
        (DATE 2025), two 2026 preprints and one GitHub study. They report aggregate scores, decision
        flips, or large models in multi-turn episodes. The paper's claim is therefore narrower: a
        syntactic-versus-semantic failure breakdown for 1--8B models across five precisions on one
        serving stack, and a full-scale test of constrained decoding under quantization. Details
        are in \texttt{notes/related\_work.md}.
\end{itemize}

% -------------------------------------------------------------------
\section{Weekly Plan}

\begin{center}
\small
\renewcommand{\arraystretch}{1.35}
\begin{tabularx}{\textwidth}{@{}c>{\raggedright\arraybackslash}p{1.7cm}LL@{}}
\toprule
\textbf{Wk} & \textbf{Dates} & \textbf{Update A (Monday)} & \textbf{Update B (Wednesday)} \\
\midrule
1 & 4--10 Oct &
  Repository set up: README, ROADMAP, UPDATES log, related-work notes, core/stretch scope agreed &
  BFCL running end-to-end on one model in FP16, with a score close to its leaderboard value \\
2 & 11--17 Oct &
  Own evaluation script that saves every raw model output and scores it with BFCL's AST checker;
  configs for all models &
  FP16 baselines for all core models \\
3 & 18--24 Oct &
  INT8 and NF4 (bitsandbytes) runs; GPTQ and AWQ 4-bit checkpoints made for every model with one
  recipe (llm-compressor or GPTQModel) &
  GPTQ and AWQ runs $\rightarrow$ \textbf{first main results table} \\
4 & 25--31 Oct &
  Failure classifier: BFCL error types sorted into syntactic (unparseable output, schema or type
  violation) and semantic (wrong function, wrong value, hallucinated or missing parameter, tool
  call on an irrelevant query) &
  \textbf{Midpoint:} first key figure (accuracy drop by precision and model size) with
  significance tests; paper outline in \texttt{paper/} \\
5 & 1--7 Nov &
  Constrained decoding (vLLM structured outputs with each tool's JSON schema) on all runs; Setup
  section drafted &
  \textbf{Headline result:} how much of the quantization penalty constrained decoding removes, and
  which semantic errors remain \\
6 & 8--14 Nov &
  Schema-stress suite v1 (about 150 template-generated cases: nested objects, enums,
  integer vs.\ float, optional fields, unicode, 50-tool menus); Introduction drafted &
  Stress-suite results broken down by failure type; Related Work section \\
7 & 15--21 Nov &
  One stretch item; latency benchmark &
  Accuracy vs.\ memory and latency curve; Method section \\
8 & 22--28 Nov &
  \textbf{Experiment freeze (22 Nov)}; all figures final; Results, Discussion and Limitations &
  Full paper draft to supervisor \\
9 & 29 Nov -- 2 Dec &
  Revisions from supervisor feedback; release preparation &
  Final paper submitted; release with reproduction steps, results and stress suite; tag
  \texttt{v1.0} \\
\bottomrule
\end{tabularx}
\end{center}

\textbf{Midpoint checkpoint (28 October).} If the baselines are not finished by the end of week 4,
drop the stretch items immediately and protect week 5 (constrained decoding) and week 8 (writing).

% -------------------------------------------------------------------
\section{Update Protocol}

\begin{itemize}
  \item \textbf{Push before every meeting}, by 12:00 noon on Monday and Wednesday, so the update is
        on GitHub before the 4:00 pm meeting. After each meeting, record the supervisor's feedback
        and action items in \texttt{notes/meetings.md}.
  \item \textbf{Commit small and often} with clear messages, so the history shows steady progress
        rather than one large push.
  \item \textbf{Keep \texttt{UPDATES.md}}, a dated log with the newest entry on top. Each entry has
        four fields: Done, Result (a number or a plot), Problem, Next. Example:
\begin{verbatim}
## 2026-10-07 (Update 2)
Done:    BFCL AST evaluation runs end-to-end on Qwen3-4B-Instruct-2507 (FP16)
Result:  simple XX.X%, multiple XX.X% (leaderboard: XX.X%)
Problem: <what blocked progress, if anything>
Next:    evaluation script with raw-output logging; all FP16 baselines
\end{verbatim}
  \item \textbf{Track the plan in GitHub Issues and Milestones}, one milestone per week, and close
        issues as they are finished.
  \item \textbf{Show something concrete in every update.} ``Read papers'' is a weak update;
        ``related-work table of 12 papers and the gap this project fills'' is a strong one.
\end{itemize}

% -------------------------------------------------------------------
\section{Repository Structure}

\begin{verbatim}
quantization-functional-call/
|-- README.md          problem, status, how to reproduce
|-- ROADMAP.md         this plan, with checkboxes
|-- UPDATES.md         dated progress log
|-- configs/           one YAML per model x precision
|-- src/
|   |-- run_eval.py    generation (vLLM), constrained decoding on/off
|   |-- score.py       BFCL AST checker wrapper
|   `-- taxonomy.py    error -> syntactic / semantic class
|-- data/stress_suite/
|-- results/           summary CSVs and figures (no model weights)
|-- notebooks/         analysis and plotting
|-- notes/             related work, design decisions, meeting notes
|-- docs/              proposal and this roadmap
`-- paper/             LaTeX source
\end{verbatim}

A \texttt{.gitignore} keeps model weights, caches and large raw outputs out of the repository.
The Hugging Face access token is never committed.

% -------------------------------------------------------------------
\section{Working Practices}

\begin{itemize}
  \item \textbf{Write the paper alongside the experiments.} Related-work notes start in week 1, the
        outline in week 4, and one section is drafted every few days from week 5, so the final
        weeks are for revision rather than first drafts.
  \item \textbf{Make every run reproducible from day one.} Use greedy decoding (temperature 0) and a
        fixed seed, and record the exact model revision and library versions in each results file.
  \item \textbf{Save raw model outputs, not only scores.} The failure taxonomy and the
        constrained-decoding analysis both depend on them, and re-running costs GPU hours.
\end{itemize}

% -------------------------------------------------------------------
\section{Key Dates}

\begin{itemize}[label=$\square$]
  \item 7 Oct --- BFCL reproduced on one model (FP16)
  \item 14 Oct --- FP16 baselines for all core models
  \item 21 Oct --- First main results table (all precisions)
  \item 28 Oct --- Midpoint: key figure and paper outline
  \item 4 Nov --- Headline result: the constrained-decoding gap
  \item 11 Nov --- Schema-stress suite results
  \item 22 Nov --- Experiment freeze
  \item 25 Nov --- Full paper draft to supervisor
  \item 2 Dec --- Final paper and \texttt{v1.0} release
\end{itemize}

\end{document}
